\begin{center}
\begin{longtblr}[
  label={tab:criterios_ie},
  caption={Criterios de inclusión y exclusión aplicados en las fases de \textit{screening} y elegibilidad, alineados con PICo, directrices de investigación (D1--D3) y criterios prácticos de Kitchenham. Elaboración propia.}]
{
    colspec = {X[m,j,3] X[m,j,5] X[m,j,5] X[m,j,1.25]},
    width=\textwidth
}
\textbf{Criterio} &
\textbf{Inclusión (Descripción operacional)} &
\textbf{Exclusión (Descripción operacional)} &
\textbf{Fase} \\

Relación con el fenómeno (Calidad del aire -- P) &
El estudio aborda explícitamente la calidad del aire o contaminación atmosférica ambiental. &
El estudio se enfoca en agua, suelo, agricultura, acuicultura, ruido u otro dominio no relacionado con calidad del aire. &
Cribado \\

Uso central de lógica difusa (I) &
El título o resumen menciona \textit{fuzzy logic} o sistemas de inferencia difusa como enfoque principal. &
No se identifica uso de lógica difusa o esta no constituye el enfoque central del estudio. &
Cribado \\

Evaluación de calidad del aire mediante lógica difusa (P + I) &
Se propone evaluación, clasificación o análisis del nivel de calidad del aire utilizando lógica difusa. &
La lógica difusa se aplica únicamente para control (PID, HVAC, secado, ventilación, hornos, etc.) sin evaluar calidad del aire como variable principal. &
Cribado \\

Naturaleza del estudio (Tipo de artículo) &
Estudio primario con propuesta técnica o implementación original. &
Revisión sistemática, revisión narrativa, \textit{mapping study}, volumen de \textit{proceedings} o literatura gris. &
Cribado \\

Plataforma o sistema de monitoreo de calidad del aire (Co) &
El estudio propone o implementa una plataforma, sistema de monitoreo, arquitectura IoT, integración de sensores o uso estructurado de \textit{datasets} para análisis de calidad del aire. &
El estudio no describe sistema, plataforma o integración de datos, o se limita a análisis teórico sin componente de monitoreo o procesamiento estructurado. &
Cribado \\

\textbf{(D1) }Especificación de variables y contaminantes utilizados &
El estudio identifica explícitamente los contaminantes atmosféricos y/o variables ambientales empleadas como entradas (p.\,ej., PM$_{2.5}$, PM$_{10}$, NO$_2$, O$_3$, CO, SO$_2$, temperatura, humedad, viento). &
No se especifican claramente las variables utilizadas o no es posible identificarlas a partir del texto completo. &
Texto Completo \\

\textbf{(D2) }Aplicación explícita de un sistema de inferencia difusa a calidad del aire &
La inferencia difusa constituye el mecanismo central para evaluar la calidad del aire y/o el riesgo asociado. &
La lógica difusa es tangencial, secundaria o no se aplica directamente a la evaluación de calidad del aire. &
Texto Completo \\

\textbf{(D2) }Descripción estructural del modelo difuso (tipo y diseño) &
Se especifica el tipo de sistema (p.\,ej., Mamdani, Sugeno, ANFIS, tipo-2 u otro) y se describen elementos estructurales (reglas, funciones de pertenencia o arquitectura). &
No se detalla el tipo de sistema difuso ni su estructura, impidiendo su caracterización y extracción sistemática. &
Texto Completo \\

\textbf{(D3) }Implementación en sistema o plataforma de monitoreo &
El estudio describe implementación en plataforma, arquitectura IoT, sistema de monitoreo o entorno estructurado de procesamiento de datos de calidad del aire. &
El modelo se presenta sin contexto de sistema/plataforma de monitoreo o sin integración operacional en un entorno de procesamiento de datos. &
Texto Completo \\

\textbf{(D3) }Descripción del entorno de aplicación y unidad de análisis (Setting / Participants) &
Se especifica el contexto de aplicación (p.\,ej., urbano/ciudad, estaciones, sensores, \textit{datasets}) y la unidad de análisis evaluada. &
No se describe claramente el entorno de aplicación o la unidad de análisis, impidiendo su comparación con otros estudios. &
Texto Completo \\

\textbf{(Transversal) }Diseño metodológico y validación documentados (Research Design / Sampling) &
El estudio describe un diseño experimental o aplicado y presenta validación empírica o evaluación cuantitativa (p.\,ej., métricas, comparación, casos de prueba). &
No presenta diseño metodológico claro y/o no reporta validación empírica o evaluación cuantitativa del sistema propuesto. &
Texto Completo \\

\end{longtblr}
\end{center}